\section{Entregables de la Parte I}

El equipo debe entregar un informe técnico con los siguientes componentes. Marque cada entregable cuando esté completo y verificable.

\begin{itemize}[leftmargin=*, topsep=2pt, itemsep=4pt, label=\checkbox]
	\item Informe técnico por equipo, con resumen ejecutivo y alcance del laboratorio.
	\item Diez fichas de caso con indicadores, análisis, hipótesis alternativas y medidas de defensa.
	\item Capturas de pantalla legibles y PCAP didácticas o sintéticas cuando corresponda, con hash SHA-256.
	\item Reglas de detección propuestas y justificadas; no se exige activar bloqueos automáticos en producción.
	\item Conclusiones sobre los límites de la evidencia y oportunidades de mejora.
	\item Anexo de evidencias con nombre de archivo, tamaño, fecha y hash SHA-256.
	\item Tabla de distribución del trabajo y lista de verificación firmada por los cuatro integrantes.
\end{itemize}

\subsection{Criterios de calidad del informe}

\begin{itemize}[leftmargin=*, topsep=2pt, itemsep=2pt]
	\item \textbf{Reproducibilidad:} cada resultado debe poder repetirse con el filtro, comando o script indicado.
	\item \textbf{Trazabilidad:} toda afirmación debe remitir a una evidencia identificada y a una marca de tiempo.
	\item \textbf{Legibilidad:} las capturas deben permitir leer direcciones, puertos, filtros y marcas de tiempo.
	\item \textbf{Anonimización:} no deben aparecer datos personales, credenciales ni direcciones de terceros.
	\item \textbf{Integridad:} la evidencia original se conserva sin modificaciones y se analiza sobre una copia de trabajo.
	\item \textbf{Postura crítica:} se deben reconocer los falsos positivos, las limitaciones y las hipótesis no confirmadas.
\end{itemize}
